\clearpage

% --------------------------------------------------------------------------------------------

\section {Summary}
\label{sec:summary}

In this technical note,
we have presented the results of a 3-flavour $\bar{\nu}_e$-appearance analysis,
performed by the {\em VALOR} group,
on the combined Run 5+6 \TotalPot{} dataset, using only data taken in antineutrino beam mode.
The goal of this analysis is to search for $\bar{\nu}_e$ appearance in a $\bar{\nu}_{\mu}$ beam.
This analysis has adapted the work presented in
Refs. \cite{OurRun1DisappAnalysis}, 
\cite{OurRun12DisappAnalysis}, 
\cite{OurRun12DisappAnalysis3Flavour},
\cite{OurRun123DisappAnalysis3Flavour},
\cite{OurRun1234aDisappAnalysis3Flavour} and 
\cite{OurRun1234DisappAnalysis3Flavour}
to this dataset. The Run 5+6 dataset corresponds to an integrated J-PARC neutrino beam
exposure of \TotalPot{} POT.
This analysis predicts \NskOscExpectation{} $\pm$ \NskOscExpectationError{} (syst)
single $e$-like ring events in Super-K, if $\bar{\nu}_e$ appearance is considered,
while \NskObserved{} were observed. 
A $\bar{\nu}_{e}$-appearance analysis was performed
in a framework of 3-flavour oscillations including matter effects in constant-density matter.
The observed reconstructed energy spectrum of single $e$-like ring events was used to 
compute the p-values.
A rate-only p-value is performed using only the total number of observed events, 
while a binned log-likelihood ratio was used to compute the rate+shape p-value,
with a test statistic, $\Delta \chi^2$, which compares the case for $\beta=0$ and $\beta=1$.
Also all \NSyst{} systematic parameters considered in this analysis were marginalized.
Among them all the oscillation parameters, except $\sin^2 \theta_{12}$ and $\Delta m^{21}$, 
as well as the mass hierarchy were marginalized.
The quoted 1-sided p-value given null $\bar{\nu}_e$ appearance, is 0.2637 for the rate-only analysis
and 0.1589 for the rate+shape analysis.
The p-value given $\bar{\nu}_e$ appearance was calculated as well and is
0.3752 for the rate-only analysis and 0.2770 for the rate+shape analysis.
the p-value given $\bar{\nu}_e$ appearance is slightly larger than the background only hypothesis. 
The fit slightly prefers the hypothesis of $\bar{\nu}_e$ appearance,
though the significance is not large enough to claim any hint of $\bar{\nu}_e$ appearance.
